\documentclass[11pt]{article}

\usepackage[margin=1.1in]{geometry}
\usepackage{amsmath,amssymb,amsthm}
\usepackage{booktabs}
\usepackage{enumitem}
\usepackage{xspace}
\usepackage{xcolor}
\usepackage[colorlinks=true,linkcolor=blue!50!black,citecolor=blue!50!black,urlcolor=blue!50!black]{hyperref}

% ---------------------------------------------------------------------------
% Macros
% ---------------------------------------------------------------------------
\newcommand{\F}{\mathbb{F}}                       % the extension field F_{2^128}
\newcommand{\Ftwo}{\F_2}
\newcommand{\eq}{\mathsf{eq}}
\newcommand{\nnz}{\mathsf{nnz}}
\newcommand{\ip}[2]{\langle #1, #2 \rangle}
\newcommand{\code}[1]{\texttt{#1}}
\newcommand{\ZC}{\textsf{Zerocheck}\xspace}
\newcommand{\LC}{\textsf{Lincheck}\xspace}
\newcommand{\pre}{\mathrm{pre}}
\newcommand{\out}{\mathrm{out}}
\newcommand{\inn}{\mathrm{in}}

\newtheorem{lemma}{Lemma}
\newtheorem{remark}{Remark}
\theoremstyle{definition}
\newtheorem{definition}{Definition}

\title{Multi-Table Proving in Flock:\\ One Proof for Heterogeneous Hash Batches\\[0.6ex]
\large Design document --- working draft v0.7}
\author{Flock}
\date{July 27, 2026}

\begin{document}
\maketitle

\begin{abstract}
Flock currently produces one proof per hash function: a proof attests to $n$
invocations of \emph{one} of BLAKE3, SHA-256, or Keccak-f[1600]. This document
designs the generalization to a \emph{single} proof attesting to a
heterogeneous batch --- e.g.\ $1000$ BLAKE3 compressions \emph{and} $4096$
SHA-256 compressions, independently invoked, with the counts chosen at prove
time --- in one proof with one commitment and one opening. The design views
the statement as a collection of \emph{tables} (one per hash type; rows =
invocations, columns = the trace of one invocation) composed into one
\emph{block-diagonal R1CS} over a \textbf{static, capacity-sized address
space}: each type owns a fixed slot sized to a maximum invocation count, the
per-proof counts are public inputs, and undeclared rows are zero --- they
satisfy the constraints identically and cost the prover essentially nothing.
The central observation is that running Flock's existing single-instance
zerocheck over the composed instance is \emph{mathematically identical} to a
batched per-table zerocheck with $\eq$-derived batching coefficients --- so
batching comes for free, with plain single-instance soundness and almost no
new zerocheck code. The genuinely new protocol work concentrates in the
lincheck, which becomes a single sumcheck over the union's column domain. The commitment layer is the jagged (stacked-columns) scheme:
only the declared rows of the used columns are committed, with the rest of
the static address space existing purely virtually. There is one protocol mode; fixed
workloads are simply the full-utilization special case.
\end{abstract}

\tableofcontents

% ===========================================================================
\section{Goals and non-goals}\label{sec:goals}

\paragraph{Goal.} One proof system execution that proves a batch containing
invocations of \emph{several different} hash functions, with the number of
invocations of each chosen at prove time. The invocations are
\textbf{independent}: no hash output feeds another hash input across (or
within) types; each row is a self-contained invocation. The multi-table
core deliberately binds \emph{no} per-invocation input/output --- all
dataflow, between tables and to the public statement, is deferred to a
planned lookup/bus layer (\S\ref{sec:statement}, \S\ref{sec:open}).

\paragraph{Design priorities} (decided up front, in order):
\begin{enumerate}[nosep]
  \item \textbf{Prover throughput.} Hashes proven per second is the headline
        metric. Per-hash witness-generation kernels and the NEON/AVX hot
        paths must not be perturbed, and the prover must pay for
        \emph{declared} work, never for capacity.
  \item \textbf{Gradual landing.} Small reviewable steps; the existing
        single-hash path stays green throughout; each phase is independently
        shippable and benchmarkable.
  \item \textbf{Static structure.} The set of hash types, their encoders,
        their widths, and their capacities come from a small predeclared
        \emph{registry}. Everything structural about the verifier --- round
        counts, control flow, precomputed digests --- depends only on the
        registry, never on the per-proof counts, which enter as scalars
        only. (This is what keeps the door open to fixed-circuit recursion
        or aggregation later.)
\end{enumerate}

\paragraph{Non-goals (for now).}
Cross-table dataflow and per-invocation public IO --- both deliberately
excluded from the multi-table core and slated for a subsequent lookup/bus
layer (\S\ref{sec:open}); a dynamic \emph{type} registry (adding a new
hash type remains a registry/code change); changing the \emph{internal}
protocol hash (the SHA-256 used for the Merkle commitment, Fiat--Shamir,
and proof-of-work grinding is untouched by this design --- ``hash'' in this
document always means the hash being \emph{proven}).

% ===========================================================================
\section{Preliminaries: the current single-table system}\label{sec:prelim}

\subsection{Notation}

$\F = \F_{2^{128}}$ is the extension field; witnesses are $\Ftwo$-valued
(bits). For $\tau, x \in \F^m$,
\[
  \eq(\tau, x) \;=\; \prod_{i=1}^{m} \bigl(\tau_i x_i + (1-\tau_i)(1-x_i)\bigr),
\]
the multilinear extension of the Kronecker delta on the Boolean cube. We use
throughout the identity
\begin{equation}\label{eq:eqsum}
  \sum_{x \in \{0,1\}^m} \eq(\tau, x) \;=\; 1
  \qquad\text{for every } \tau \in \F^m,
\end{equation}
which holds because each factor sums to $\tau_i + (1-\tau_i) = 1$ over
$x_i \in \{0,1\}$. For a vector $v \in \F^{2^m}$, $\hat v$ denotes its
multilinear extension (MLE); index bits are little-endian (bit $0$ is the
least significant). For a point $r \in \F^m$ and a bit range $[a,b)$ we write
$r[a{:}b]$ for the corresponding subvector of coordinates.

\subsection{The single-table instance}\label{sec:single}

A Flock statement today is a \emph{block R1CS}
(\code{flock-core/src/r1cs.rs}, \code{struct BlockR1cs}): one \emph{base
block} of sparse $\Ftwo$-matrices $(A_0, B_0, C_0)$, each $k \times k$ with
$k = 2^{\kappa}$, with $C_0 = I$ (``circuit form'': the $C$-side just reads
the witness), replicated $n = 2^{\nu}$ times. The full constraint matrix is
implicitly
\[
  A \;=\; I_{n} \otimes A_0, \qquad B \;=\; I_{n} \otimes B_0,
  \qquad C \;=\; I_{n} \otimes C_0 = I,
\]
and the witness $z \in \Ftwo^{2^{m}}$, $m = \nu + \kappa$, satisfies
$(Az) \circ (Bz) = Cz$ (Hadamard product), i.e.\ each block independently
satisfies the base relation. Only the first $u \le k$ columns of a block are
``useful''; columns $[u, k)$ are zero padding, recorded by a single scalar
\code{useful\_bits}.

\paragraph{The table view.} This is exactly a \emph{table}: $n$ rows (one per
hash invocation), $k$ columns (the trace of one invocation --- for the current
encoders, $\kappa = 14$ for BLAKE3, $15$ for SHA-256, $16$ for Keccak, $17$
for the 3-wide Keccak variant; all share the zerocheck skip parameter
$k_{\mathrm{skip}} = 6$). The witness index splits as
\[
  x \;=\; (\underbrace{x[0{:}\kappa]}_{\text{column / ``inner''}},\;
           \underbrace{x[\kappa{:}m]}_{\text{row / ``outer''}}).
\]
This is today's default \emph{row-major} convention (one invocation's
trace is contiguous); the multi-table design adopts the code's other,
chunk-column-major convention (\code{WitnessLayout::BatchMajor},
\S\ref{sec:union}).

\subsection{The PIOP, briefly}\label{sec:pioponetable}

Let $a = Az$, $b = Bz$, $c = Cz = z$ over the $2^m$-point domain. The proof
runs:

\paragraph{\ZC (\code{zerocheck.rs}).} Proves $\hat a(x)\hat b(x) = \hat c(x)$
for all Boolean $x$: the verifier samples $\tau \in \F^{m}$, the prover runs a
sumcheck for
\begin{equation}\label{eq:zc-single}
  \sum_{x \in \{0,1\}^{m}} \eq(\tau, x)\,
     \bigl(\hat a(x)\hat b(x) - \hat c(x)\bigr) \;=\; 0 .
\end{equation}
The low $k_{\mathrm{skip}} = 6$ variables are bound in a single
\emph{univariate-skip} round (one univariate over a $2^{6}$-point domain),
which is where most of the prover's win from $\Ftwo$-valued data lives; the
remaining $m - 6$ rounds are ordinary. The output is a random point
$r \in \F^{m}$ and claims $\hat a(r), \hat b(r), \hat c(r)$. Since $C = I$,
$\hat c(r) = \hat z(r)$ is already a witness evaluation claim. A
\code{PaddingSpec} $= (\kappa, u)$ lets the prover skip work on the uniform
zero-padding columns.

\paragraph{\LC (\code{lincheck.rs}).} Reduces $\hat a(r)$ and $\hat b(r)$
(batched by a fresh $\alpha$: the claim is $\alpha\hat a(r) + \hat b(r)$) to a
witness evaluation. The load-bearing structure is the factorization of the
replicated matrix's MLE:
\begin{equation}\label{eq:factor-single}
  \hat A\bigl((r_{\out}, r_{\inn}), (c_{\out}, c_{\inn})\bigr)
  \;=\; \eq(r_{\out}, c_{\out}) \cdot \hat A_0(r_{\inn}, c_{\inn}),
\end{equation}
so that
\[
  \hat a(r) = \sum_{c} \hat A(r, c)\,\hat z(c)
  = \sum_{c_{\inn} \in \{0,1\}^{\kappa}}
      \hat A_0(r_{\inn}, c_{\inn})\;
      \underbrace{\hat z(c_{\inn}, r_{\out})}_{\text{partial eval.\ of } \hat z},
\]
where the sum over $c_{\out}$ \emph{collapsed} into a partial evaluation of
the witness MLE at the row point $r_{\out}$. The prover materializes
$y(c_{\inn}) := \hat z(c_{\inn}, r_{\out})$ (one $O(2^m)$ fold of the witness)
and the \emph{comb} vector
$\mathrm{comb}(c_{\inn}) = \alpha\,(\eq_{r_{\inn}}^{\top} A_0)(c_{\inn})
 + (\eq_{r_{\inn}}^{\top} B_0)(c_{\inn})$,
then runs a $\kappa$-bit sumcheck for
$\sum_{c_{\inn}} \mathrm{comb}(c_{\inn})\, y(c_{\inn})$. The verifier's matrix
work is $O(\nnz(A_0) + \nnz(B_0))$, \textbf{independent of the row count} ---
it never touches the expanded matrix. (In code the comb evaluation goes
through \code{trait LincheckCircuit}; Keccak implements it with a custom
walker rather than literal sparse matrices. Everything below interacts with
that trait, not with matrices directly.)

\paragraph{Claims, glue, opening.} All evaluation claims on $\hat z$ --- from
\ZC ($\hat c(r)$), from \LC, and from the statement-binding arguments (e.g.\
today's chain argument, which constrains input/output \emph{regions} of the
committed witness via an auxiliary shift-sumcheck; note it is claim-level
``glue'', \emph{not} extra R1CS rows) --- are carried to the committed
polynomial by a two-stage machinery: \emph{ring-switching}
(\code{pcs/ring\_switch.rs}) eliminates the seven in-word packing
coordinates --- including the univariate skip coordinate, which is why
downstream layers never see quirky structure --- turning each bit-level
claim into an inner-product claim, against a transparent
(verifier-evaluable) weight, over the \emph{packed} ($\F$-word)
polynomial; the claims are then $\gamma$-batched
(\code{pcs::open\_batch\_mixed\ldots}, label
\code{flock-pcs-open-batch-v0}) into one Ligerito opening of the single
committed witness polynomial.

The design below preserves every element of this pipeline and generalizes the
instance it runs on.

% ===========================================================================
\section{The multi-table model}\label{sec:model}

\begin{definition}[Type registry]\label{def:registry}
The \emph{registry} fixes, once, an ordered list of \emph{table types}
$t = 1, \dots, T$. Type $t$ consists of: a base block
$(A_t, B_t, C_t{=}I)$ of width $k_t = 2^{\kappa_t}$ with $u_t \le k_t$
useful columns, and a \emph{row capacity} $2^{\nu_t^{\max}}$ --- the padded
maximum: the largest number of invocations of type $t$ a single proof may
declare, sizing the type's (virtual) slot once and for all. The number of
rows \emph{actually used} is a per-proof quantity, the count $n_t$ of
Definition~\ref{def:instance}; the prover pays only for those. Write
$m_t = \nu_t^{\max} + \kappa_t$; the type's \emph{capacity area} is
$s_t = 2^{m_t}$. The registry orders types by non-increasing capacity area
$s_1 \ge s_2 \ge \cdots \ge s_T$.
\end{definition}

Definition~\ref{def:registry} permits per-type capacities, but the
\textbf{adopted convention is uniform row capacity}:
$\nu_t^{\max} \equiv \nu$ for all $t$. Uniformity is load-bearing
downstream: every slot then shares the same row coordinates (one shared
row-fold table and a \emph{single} output claim in the lincheck,
\S\ref{sec:lc}), and the packed union polynomial becomes literally the
jagged grid polynomial, so union-level claims transport without per-slot
decomposition (\S\ref{sec:commit}). The cost is only virtual: a
rarely-used wide type is granted the same row capacity as the workhorse
types, but undeclared capacity is address space, not sumcheck work
(\S\ref{sec:zc-sparse}) --- nor commit work: under height-$n_t$ stacking
the committed size is count-proportional (\S\ref{sec:commit}) --- and the
round count exceeds that of an
area-balanced registry ($\nu_t^{\max} + \kappa_t$ constant) by at most
$\kappa_{\max} - \kappa_{\min}$ ($\le 3$ today). Non-uniform capacities
remain expressible if address-space pressure ever demands them, at the
price of re-introducing per-slot claim decomposition.

Two clarifications on Definition~\ref{def:registry}. First, capacities are
powers of two \emph{by fiat}: slots must be subcubes of the global cube
(Lemma~\ref{lem:align} and everything built on it), and subcubes have
power-of-two size. This costs nothing --- capacity is virtual, so a true
maximum of, say, $3 \cdot 2^{10}$ is simply rounded up to $2^{12}$; the
never-used rows are skipped by the prover, and all genuine arbitrariness
lives in the counts $n_t$, which are unrestricted. Second, the registry
ordering is by capacity \emph{area} $s_t$ (rows $\times$ width), not by row
capacity alone --- a wider type with fewer maximum rows may sort first ---
and the ordering is itself only a device: it is the canonical way to obtain
the alignment property of Lemma~\ref{lem:align}; under the
uniform-capacity convention it is simply a sort by width.

\begin{definition}[Instance]\label{def:instance}
A proof \emph{instance} consists of public \emph{counts}
$(n_t)_{t=1..T}$ with $0 \le n_t \le 2^{\nu_t^{\max}}$ --- arbitrary
integers, \emph{not} necessarily powers of two --- together with, per type
$t$, the witness table $Z_t \in \Ftwo^{n_t \times k_t}$: row $r < n_t$ holds
the full trace of invocation $r$ and independently satisfies
$(A_t z) \circ (B_t z) = z$. Rows $[n_t, 2^{\nu_t^{\max}})$ are \emph{dummy
rows}, identically zero.
\end{definition}

Independence of invocations means there are no constraints connecting
different rows, neither within a table nor across tables. The multi-table
layer attaches no statement to individual rows --- no input/output binding
(\S\ref{sec:statement}); in particular, dummy rows are referenced by
nothing at all.

If a workload exceeds a capacity, the proof is split (or a larger registry
tier is used, \S\ref{sec:statement}); capacities are a registry-design
question, not a protocol limit (\S\ref{sec:open}).

% ===========================================================================
\section{The union instance}\label{sec:union}

We compose the tables into one index space and one R1CS. The layout is
computed \textbf{once, at registry-definition time}, from the capacity
areas --- the per-proof counts never move anything.

\subsection{Static layout and alignment}

\paragraph{Layout convention (BatchMajor).} Within a slot, addresses use
the \emph{implemented} \code{WitnessLayout::BatchMajor} order
(\code{r1cs.rs}): trace variables are grouped into chunks of $128$; one
$128$-bit word holds one invocation's chunk; the words of a chunk are
contiguous across invocations. The address split is therefore
\emph{interleaved}:
\[
  \text{addr} = [\,\underbrace{7 \text{ in-word bits}}_{\text{column,
  low}} \mid \underbrace{\nu_t^{\max} \text{ row bits}}_{\text{invocation}}
  \mid \underbrace{\kappa_t - 7 \text{ chunk bits}}_{\text{column,
  high}}\,],
\]
i.e.\ slot $t$'s \emph{row coordinates} are $R_t = [7,\, 7{+}\nu_t^{\max})$
and its \emph{column coordinates} are the non-contiguous
$C_t = [0,7) \cup [7{+}\nu_t^{\max},\, m_t)$. Under the uniform-capacity
convention (\S\ref{sec:model}), $R_t = [7, 7{+}\nu) =: R$ is the
\emph{same} coordinate set for every slot; we write $R$ where uniformity
matters. The interleaving is
harmless: all the math below depends only on the \emph{partition} of a
slot's coordinates into (row, column, prefix) groups, never on their
contiguity. The convention earns its keep three ways: the univariate skip
runs over the low six in-word (column) bits \emph{exactly as today} ---
the existing skip kernels are byte-level layout-agnostic and run unchanged
(\S\ref{sec:skip}); chunk-columns are contiguous across invocations, the
stacking order the jagged commitment needs (\S\ref{sec:commit}), at
$128$-variable-chunk granularity; and the witness drivers already emit
this layout (with per-block padding coalesced), producing $z$, $Az$, and
$Bz$ during generation --- no matrix application on the hot path. Requires
$\kappa_t \ge 7$ (all current $\kappa_t \ge 14$).

Let $M$ be minimal with $2^{M} \ge \sum_t s_t$ (choose capacities so that
$\sum_t s_t = 2^{M}$ exactly if desired; a shortfall is harmless by
Remark~\ref{rem:padding}). Assign type $t$ the \emph{slot}
$[o_t, o_t + s_t)$ of $\{0,1\}^{M}$, where $o_t = \sum_{t' < t} s_{t'}$.

\begin{lemma}[Alignment]\label{lem:align}
With the descending-area ordering, $o_t$ is a multiple of $s_t = 2^{m_t}$.
Consequently, for a global index $x$ in slot $t$:
\[
  x|_{R_t} = \text{row index},\quad
  x|_{C_t} = \text{column index},\quad
  x[m_t{:}M] = p_t \;(\text{a fixed bit pattern}, \; p_t = o_t / 2^{m_t}),
\]
where $x|_S$ denotes the restriction of $x$ to the coordinate set $S$ and
$(R_t, C_t)$ is the interleaved row/column partition of the BatchMajor
convention above.
\end{lemma}
\begin{proof}
$o_t$ is a sum of terms $2^{m_{t'}}$ with $m_{t'} \ge m_t$, each divisible by
$2^{m_t}$.
\end{proof}

We call $p_t$ the \emph{slot prefix}: the Boolean string of length
$M - m_t$ that the top bits of \emph{every} address in slot $t$ are frozen
to (numerically $o_t / 2^{m_t}$, an exact division by the lemma). Note its
length varies per slot --- larger slots have shorter prefixes --- so
expressions like $\eq(\tau[m_t{:}M], p_t)$ always pair it with a
coordinate range of matching length. Example ($M = 6$, slot sizes
$32, 16, 16$): slots $[0,32)$, $[32,48)$, $[48,64)$ have prefixes
$p_1 = \mathtt{0}$, $p_2 = \mathtt{10}$, $p_3 = \mathtt{11}$.

So each slot is a subcube of $\{0,1\}^M$, selected by fixing the top
$M - m_t$ bits to the static prefix $p_t$, and \emph{within} the
subcube the index has the same (column, row) split as a single-table
instance. Distinct slots are disjoint subcubes. (The alignment property is
what all the math below needs; sorting is merely how the registry obtains
it once and for all.)

\subsection{The block-diagonal R1CS}

The union instance has constraint matrices
\[
  A = \mathrm{diag}\bigl(I_{2^{\nu_1^{\max}}} \otimes A_1,\; \dots,\;
                         I_{2^{\nu_T^{\max}}} \otimes A_T\bigr),
  \qquad B \text{ likewise}, \qquad C = I_{2^M},
\]
and witness $z \in \Ftwo^{2^M}$ = the concatenation of the per-type
witnesses in slot order (each table internally BatchMajor, per the
convention of \S\ref{sec:union}), with dummy rows zero. The satisfied relation is, as before, pointwise:
$(Az) \circ (Bz) = z$ over all of $\{0,1\}^M$.

\begin{remark}[Dummy rows are complete; padding is self-enforcing]
\label{rem:padding}
The R1CS is homogeneous --- constants enter only through the constant-pin
column --- so the all-zero row satisfies
$(A_t \cdot 0)\circ(B_t \cdot 0) = 0 = C_t \cdot 0$, with the constant pin
at $0$ marking the row as dummy; no witness-generation work is performed
for dummy rows. Soundness needs nothing from them either: nothing in the
system references dummy rows (the core binds no per-row IO at all,
\S\ref{sec:statement}). Under the jagged commitment (\S\ref{sec:commit})
dummy rows are \emph{dropped} from the committed stack, along with the
useless chunks and the inter-slot gaps (height-$n_t$ stacking); every
dropped word of the padded virtual buffer is zero --- the partial witness
drivers zero dummy rows (pin included) and never write useless columns or
gaps, an invariant the transport debug-asserts --- and the declared rows
are marked by the lincheck's count-derived pin target
(Remark~\ref{rem:pinbind}). In any gap
region beyond the last slot, $A$ and $B$ have zero rows
and $C = I$, so the constraint reads $0 \cdot 0 = z(x)$, \emph{forcing}
zero. Padding therefore needs no soundness caveats; its only cost is
prover work, which \S\ref{sec:zc}--\ref{sec:lc} reduce to zero.
\end{remark}

There are \textbf{no selector polynomials} anywhere: disjointness of the
slots is the selection mechanism. The ``selector'' will surface only as
scalar $\eq$-prefix weights in the verifier's algebra, essentially for free.

% ===========================================================================
\section{Zerocheck over the union}\label{sec:zc}

Run \emph{exactly the existing zerocheck} (\S\ref{sec:pioponetable}) on the
union instance: sample $\tau \in \F^{M}$, prove
\begin{equation}\label{eq:zc-union}
  \sum_{x \in \{0,1\}^{M}} \eq(\tau, x)
     \bigl(\hat a(x)\hat b(x) - \hat c(x)\bigr) = 0 ,
\end{equation}
by an $M$-round sumcheck ($1$ univariate-skip round $+$ $M - 6$ ordinary
rounds), where now $a = Az$ etc.\ for the block-diagonal $A$. Note $M$ is
registry-determined, hence \emph{static}: a few more rounds than a
perfectly-sized instance, in exchange for a count-independent verifier.

\subsection{Why this is implicit batching}

Split the sum in \eqref{eq:zc-union} by slot. For $x$ in slot $t$,
$\eq(\tau, x)$ factors across the bit split of Lemma~\ref{lem:align}:
\[
  \eq(\tau, x) \;=\;
  \underbrace{\eq\bigl(\tau[m_t{:}M],\, p_t\bigr)}_{=: \, w_t \;\in\; \F}
  \;\cdot\;
  \eq\bigl(\tau[0{:}m_t],\, x[0{:}m_t]\bigr),
\]
and the constraint polynomial restricted to slot $t$ is exactly type
$t$'s own $\hat a_t \hat b_t - \hat c_t$. Hence
\begin{equation}\label{eq:implicit-batch}
  \eqref{eq:zc-union}
  \iff
  \sum_{t=1}^{T} w_t \cdot
  \underbrace{\sum_{x' \in \{0,1\}^{m_t}}
     \eq\bigl(\tau[0{:}m_t], x'\bigr)
     \bigl(\hat a_t \hat b_t - \hat c_t\bigr)(x')}_{\text{type } t\text{'s
     single-table zerocheck sum}} \;=\; 0 .
\end{equation}
That is: the union zerocheck \emph{is} a batched per-table zerocheck, with
batching coefficients $w_t = \eq(\tau[m_t{:}M], p_t)$ derived from the shared
randomness $\tau$ instead of independently sampled. There is no separate
batching protocol, no batching soundness lemma (soundness is the ordinary
single-instance zerocheck argument over $\{0,1\}^M$: if any invocation of any
table is wrong, the constraint polynomial is nonzero somewhere on the cube,
and \eqref{eq:zc-union} fails except with probability
$O(M/|\F|)$ over $\tau$ plus the sumcheck error), and no scheduling
machinery. Different tables simply live at different (static) addresses.

\subsection{The prover pays for declared work only}\label{sec:zc-sparse}

On every region where the witness --- hence $a$, $b$, and $c$ --- is
identically zero (dummy rows, useless columns, the gap), the summand of
\eqref{eq:zc-union} vanishes identically and contributes nothing to any
round polynomial. The \code{PaddingSpec} generalizes from today's single
$(\kappa, u)$ pair to a \textbf{run-list derived from the public counts}:
per slot, columns $[u_t, k_t)$ and rows $[n_t, 2^{\nu_t^{\max}})$ are dead
(this generalizes \code{useful\_bits} skipping from columns to rows).
Concretely, under BatchMajor the run-list's blocks are
\emph{chunk-columns} ($k\_log = 7 + \nu$): slot $t$ contributes
$\lceil u_t/128 \rceil$ used chunk-column blocks whose useful prefix is
$128\, n_t$ bits (dummy rows are the dead suffix of every chunk-column),
followed by an \emph{explicit} zero run for its useless chunks ---
mid-list gaps must be explicit so later slots keep their offsets --- with
only the region past the last slot left to the implicit trailing gap.
With run-list--driven iteration the prover's sumcheck work is
proportional to
$\sum_t n_t k_t$ --- the \emph{declared} area --- while the round count
stays the static $M$.

\subsection{The univariate skip survives globally --- unchanged}
\label{sec:skip}

The skip merges the six low variables of the sumcheck into one univariate
round; its prover win --- the six most expensive rounds computed directly
on $\Ftwo$-valued data --- requires only that each merged group of $64$
domain points be memory-contiguous. Under the BatchMajor convention the
global low six bits of any index are in-word \emph{column} bits for every
slot (all $\kappa_t \ge 7$, and every slot offset is a multiple of
$2^{m_t}$), exactly as in today's row-major layout: a skip group is $64$
adjacent trace variables of one invocation, one $u64$ of a $128$-bit
word. The existing round-1/round-2 skip kernels are byte-level
layout-agnostic and run \emph{unchanged} --- no new kernels, no
re-orientation. The zerocheck constraint is pointwise, so mixing slots
changes nothing else: one skip round, all slots at once, no per-table
scheduling. The skip's univariate coordinate sits among the \emph{column}
coordinates of the output point, precisely as today, so the lincheck's
quirky-point handling (\S\ref{sec:lc}) and the existing claim-point
conventions carry over; the coordinate is then eliminated by the existing
ring-switching stage before the jagged transport ever sees a point
(\S\ref{sec:commit}).

\subsection{What changes in code}

Only bookkeeping:
\begin{itemize}[nosep]
  \item \code{BlockR1cs} generalizes to the registry schedule of slots
        $\{(\text{base block}, \kappa_t, u_t, \nu_t^{\max})\}$ (one entry
        $\equiv$ today's struct); \code{apply\_a} applies each base block
        over its slot's declared rows.
  \item \code{PaddingSpec} $(\kappa, u)$ becomes the count-derived run-list
        of \S\ref{sec:zc-sparse}.
  \item The zerocheck sumcheck kernels: \textbf{unchanged, byte-identical}
        (pointwise and layout-agnostic; the code map confirmed the uniform
        \code{PaddingSpec} was the only single-table assumption, and the
        BatchMajor prover paths for BLAKE3/SHA-256 already exercise these
        kernels today).
\end{itemize}

\paragraph{Output.} A point $r \in \F^{M}$ and three claims
$\hat a(r), \hat b(r), \hat c(r)$, with $\hat c(r) = \hat z(r)$ a direct
witness claim.

% ===========================================================================
\section{Lincheck over the union}\label{sec:lc}

This is where the real (and modest) protocol work lives. The zerocheck left
two matrix claims, $\hat a(r)$ and $\hat b(r)$. For clarity we develop the
lincheck for the $A$-claim alone, $v_A := \hat a(r)$;
\S\ref{sec:batchAB} then observes that the $B$-claim is word-for-word
identical and that the two sumchecks merge into one at the cost of a single
random scalar.

\subsection{Closed form of the union matrix MLE}

Combining the block-diagonal structure with
\eqref{eq:factor-single} per slot and Lemma~\ref{lem:align}:
\begin{equation}\label{eq:factor-union}
  \hat A(r, c) \;=\; \sum_{t=1}^{T}
     \eq\bigl(r[m_t{:}M], p_t\bigr)\,
     \eq\bigl(c[m_t{:}M], p_t\bigr)\,
     \eq\bigl(r|_{R_t},\, c|_{R_t}\bigr)\,
     \hat A_t\bigl(r|_{C_t},\, c|_{C_t}\bigr),
\end{equation}
\paragraph{Derivation of \eqref{eq:factor-union}.} By definition
$\hat A(r,c) = \sum_{x,y \in \{0,1\}^M} \eq(r,x)\,\eq(c,y)\,A[x,y]$, and
the MLE is linear in the matrix, so split $A$ by slot; cross-slot entries
vanish, and slot $t$'s part is supported on pairs with both prefixes frozen
at $p_t$. Since $\eq$ is a product over coordinates, it factors over any
partition of them: for $x$ in slot $t$,
$\eq(r,x) = \eq(r[m_t{:}M], p_t)\cdot\eq(r|_{R_t}, x|_{R_t})
\cdot\eq(r|_{C_t}, x|_{C_t})$ --- $\eq$ factors over \emph{any} partition
of the coordinates, contiguous or not --- and the prefix factor ---
constant over the slot --- pulls out of the sum (likewise for $c$): these
are the first two factors of \eqref{eq:factor-union}, the multilinear
indicators ``$r$ addresses slot $t$'' and ``$c$ addresses slot $t$''. What
remains is the MLE of $I \otimes A_t$ in the within-slot coordinates, and
the MLE of a Kronecker product is the product of the factors' MLEs: the
identity contributes $\eq(r|_{R_t}, c|_{R_t})$
(``same invocation'') via
\begin{equation*}
  \sum_{w \in \{0,1\}^{\nu}} \eq(a, w)\,\eq(b, w) \;=\; \eq(a, b),
\end{equation*}
and the base block contributes $\hat A_t$. Note each term partitions the
\emph{same} shared coordinates by its own $(R_t, C_t, \text{prefix})$ ---
legal, since every term is simply a function of all $M$ coordinates.
Sanity check: at a Boolean pair $(x, y)$, every term except the one for
the slot actually containing $x$ dies at its prefix factor, and the
surviving term reproduces the matrix entry.

The verifier can evaluate the scalar prefix weights in $O(T \cdot M)$
additions and multiplications; the per-type base-block evaluations cost
$O(\sum_t \nnz(A_t, B_t))$ via the existing \code{LincheckCircuit} comb
machinery --- still never touching any expanded matrix.

\subsection{Per-slot collapse}

Substituting \eqref{eq:factor-union} into
$v_A = \sum_{c \in \{0,1\}^M} \hat A(r, c)\, \hat z(c)$ and
collapsing, per slot, the sum over the row coordinates $c|_{R_t}$ into a
partial evaluation of that slot's witness MLE (exactly as in the
single-table case):
\begin{equation}\label{eq:lc-claim}
  v_A \;=\; \sum_{t=1}^{T} w_t(r) \sum_{c' \in \{0,1\}^{\kappa_t}}
      \mathrm{comb}^A_t(c')\; y_t(c'),
  \qquad
  \begin{aligned}
    w_t(r) &= \eq\bigl(r[m_t{:}M], p_t\bigr),\\
    y_t(c') &= \hat z_t\bigl(r|_{R},\, c'\bigr),
  \end{aligned}
\end{equation}
where
$\mathrm{comb}^A_t(c') = \sum_{x'} \eq(r|_{C_t}, x')\, A_t[x', c']$
is the base block folded by the weight table of the within-block
constraint coordinates of $r$ --- the \emph{quirky} table, since $C_t$
contains the univariate skip coordinate (\S\ref{sec:skip}), computed
through the type's \code{LincheckCircuit} \emph{exactly as in today's
single-table code} (the code's ``semantic \code{QuirkyPoint}'' is
layout-independent) --- and $\hat z_t(\cdot,\cdot)$ is the MLE of slot
$t$'s (zero-padded) witness, rows-coordinates first. The prover
materializes each $y_t$ by folding slot $t$'s \emph{declared rows} with
the standard $\eq$ table of the \emph{shared} row coordinates $r|_{R}$
--- one table serves every slot, by uniform capacity ---
$O(\sum_t n_t k_t)$ in total, the declared area, analogous to today's
single outer fold. The division of labor (quirky on the comb side,
standard $\eq$ on the row-fold side) is \emph{identical} to today's
single-table lincheck.

Note the weights $w_t(r)$ are \emph{fixed known scalars} (determined by $r$
and the registry), not fresh randomness: \eqref{eq:lc-claim} is a single
claim about a fixed weighted sum, so no batching randomness is needed
\emph{across slots} (contrast \S\ref{sec:batchAB}, where combining the
$A$- and $B$-claims genuinely does use a random $\alpha$).

\subsection{One sumcheck over the union column domain}\label{sec:multisize}

What remains is to prove \eqref{eq:lc-claim}: a sum of $T$ sub-sums over
column domains of different sizes $2^{\kappa_t}$. Do not batch them ---
\emph{address} them, exactly as the zerocheck did. Under the
uniform-capacity convention all slots share the row coordinates $R$, so
the per-slot folded vectors are restrictions of \emph{one} function
\[
  g(c) \;:=\; \hat z\bigl(\text{rows} \leftarrow r|_R\bigr)(c),
  \qquad c \in \{0,1\}^{M-\nu},
\]
the union with its row block partially evaluated, with $y_t = g$
restricted to slot $t$'s column subcube (the slot's in-word and chunk
coordinates, its column prefix fixed; areas $2^{\kappa_t}$, dyadic, and
aligned by the argument of Lemma~\ref{lem:align} applied to the column
space). Absorb the fixed weights into the public side by defining
\[
  \mathrm{Comb}(c) \;:=\; w_t(r)\cdot \mathrm{comb}^A_t(c')
  \;\;\text{on slot } t\text{'s subcube}, \qquad 0 \;\text{on gaps}.
\]
Then \eqref{eq:lc-claim} reads
\begin{equation}\label{eq:union-lc}
  v_A \;=\; \sum_{c \in \{0,1\}^{M-\nu}} \mathrm{Comb}(c)\, g(c):
\end{equation}
\emph{one} ordinary product sumcheck over
$M - \nu \approx \kappa_{\max} + \log T$ variables. Slots that have
collapsed contribute $O(1)$ scalar work in the remaining rounds (the same
implementation pattern as the zerocheck fold, \S\ref{sec:open}). At the
end the verifier evaluates $\widehat{\mathrm{Comb}}(c^*)$ by the usual
closed form --- per-type comb machinery at $O(\nnz)$ per type, times
prefix-$\eq$ weights --- and is left with \textbf{a single witness claim}
\[
  \hat g(c^*) \;=\; \hat z\bigl(c^*\text{-columns},\; r|_R\text{-rows}\bigr),
\]
structurally identical to today's single-table lincheck output.

\begin{remark}[Why not batch per-slot sumchecks]\label{rem:char2}
An earlier draft batched $T$ per-slot sumchecks of different sizes by a
late-join mechanism, which in characteristic 2 requires an $\eq$-weighted
embedding (the textbook $2^{-\Delta}$ rescaling divides by two, and
$2 = 0$ in $\F_{2^{128}}$). The union-domain formulation supersedes all of
it: same prover work, about $\log T$ more rounds, no auxiliary embedding
point, and one output claim instead of $T$. The lesson generalizes:
wherever this design meets sums over different-sized cubes, aligned
addressing beats batching machinery.
\end{remark}

\subsection{The \texorpdfstring{$B$}{B}-claim, and batching the two
sumchecks}\label{sec:batchAB}

Everything above applies verbatim to the $B$-claim $v_B := \hat b(r)$, with
$\mathrm{comb}^B_t$ folding $B_t$ instead of $A_t$. Crucially, the folded
witness vectors $y_t$ \emph{do not change} --- they depend only on $r$ and
the witness, never on the matrix. The two multi-size sumchecks therefore
range over identical domains against identical $y_t$'s and differ only in
the comb factor, so they merge into one: squeeze a fresh $\alpha$ from the
transcript (after both claims are bound) and run the sumcheck of
\S\ref{sec:multisize} once for the combined claim
\[
  \alpha\, v_A + v_B \;=\; \sum_{t=1}^{T} w_t(r) \sum_{c'}
    \bigl(\alpha\,\mathrm{comb}^A_t + \mathrm{comb}^B_t\bigr)(c')\;
    y_t(c') .
\]
This is standard random-linear-combination batching of two claims (error
$1/|\F|$), costs one scalar and zero extra rounds, and the combined comb
$\alpha\,\mathrm{comb}^A_t + \mathrm{comb}^B_t$ is exactly what the
existing \code{fold\_alpha\_batched} computes --- the code's
$\alpha$-comb is this batching, seen from the single-table special case.
The end-of-sumcheck outputs --- per-type comb evaluations and the witness
claims $y_t(c^*_t)$ --- are shared between the two claims, so batching
adds nothing there either.

\begin{remark}[The counts bind inside the lincheck]\label{rem:pinbind}
The implementation binds the declared counts a \emph{second} time, inside
the PIOP itself. Each type's constant-pin column joins the comb with a
fresh transcript scalar $\beta_t$ (exactly as in today's single-table
lincheck), and the verifier's consistency target gains the
\emph{count-derived} term
$\beta_t \cdot \sum_{\mathrm{row} < n_t} \eq(r|_R, \mathrm{row})$:
declared rows carry the pin at $1$ and dummy rows are all-zero
\emph{including the pin} (Remark~\ref{rem:padding}), so this partial
$\eq$-sum --- computable in $O(\nu)$ field operations --- is exactly what
the honest pin column folds to. At full utilization the sum is $1$ by
\eqref{eq:eqsum}, reproducing the single-table target verbatim. A wrong
declared count is therefore rejected by the lincheck itself, independently
of the transcript binding (\S\ref{sec:statement}) and of the commitment
layout (\S\ref{sec:commit}) --- defense in depth the original design did
not require.
\end{remark}

\subsection{Verifier cost summary}

Rounds: $M$ (zerocheck) $+$ $M - \nu$ (lincheck) --- registry-static,
independent of the counts, growing only logarithmically in $T$. Claims:
$3$ zerocheck scalars and \emph{one} lincheck witness claim --- no
per-slot sub-claims anywhere. Matrix work: $O(\sum_t \nnz(A_t, B_t))$,
linear in the number of types --- the one term that grows if the registry
grows large; see \S\ref{sec:open} for the delegation escape hatch. Prefix
weights: $O(T \cdot M)$. Weight structure, comb structures, and the
registry digest are all precomputable; the counts enter only through run
lists and fold extents. Under the uniform-capacity convention the
zerocheck's $\hat c(r) = \hat z(r)$ claim needs no per-slot decomposition
either: it transports directly (\S\ref{sec:commit}).

% ===========================================================================
\section{The commitment layer: jagged transport}\label{sec:commit}

The PIOP above emits exactly two evaluation claims on the union
polynomial --- the zerocheck's $\hat c(r) = \hat z(r)$ and the lincheck's
$\hat g(c^*)$ --- plus, on single-type statements, the existing glue
claims (a future lookup/bus layer will add its own, \S\ref{sec:open}).
The commitment layer
turns these into openings of \emph{one} committed polynomial while
committing \textbf{only the declared rows of the used chunk-columns}: the
dead part of the capacity-sized address space of \S\ref{sec:union} ---
dummy rows included --- exists purely virtually (Reduction 2 below). The
pipeline is best
understood as \textbf{three polynomials and two reductions}:
\begin{enumerate}[nosep]
  \item the \emph{sparse bit polynomial} $\hat z$ ($M$ variables): the
        arithmetization's object --- bit-valued, zero on dummy rows,
        useless columns, and gaps. The zerocheck and lincheck speak only
        about this polynomial.
  \item the \emph{sparse packed polynomial} $\hat f$ ($M-7$ variables):
        each $128$ bits of $\hat z$ packed into one $\F$-word, still over
        the virtual word-address space, still zero on dead words. Never
        committed, never materialized beyond its support.
  \item the \emph{dense packed polynomial} $\hat q$
        ($M_{\mathrm{dense}}$ variables): the sparse packed polynomial
        with the dead chunk-columns deleted --- the only polynomial
        Ligerito commits and opens.
\end{enumerate}
Ring-switching reduces claims about (1) to a claim about (2); the jagged
transport reduces a claim about (2) to a claim about (3). Each reduction
discharges exactly one mismatch --- bit-versus-word, then
virtual-versus-dense. The two reductions use different \emph{claim
types}, which dictates one additional, standard step between them.

\paragraph{Reduction 1: ring-switching (existing, unchanged).} A PIOP
claim is bit-level and carries the univariate skip coordinate among its
low (in-word) coordinates. Ring-switching sends the $128$-element
slice-fold $\hat s_v$; the verifier checks the bit-level claim against it
using the in-word weights --- this is where the skip coordinate is
consumed, never to reappear downstream. Binding the prover-supplied
$\hat s_v$ to the commitment then \emph{requires} fresh randomness (a
single opening would pin down only one linear functional of a
$128$-dimensional message), so the output claim about $\hat f$ is
inherently an \textbf{inner product} $\langle f, W\rangle = t$ against a
transparent, verifier-evaluable weight $W$ --- and $W$ does \emph{not}
factor as the $\eq$ tensor of any point, because its per-entry bit
surgery does not commute with field multiplication. Today's machinery
already $\gamma$-batches \emph{all} claims (ab, c, glue) into one such
statement $(b_{\mathrm{combined}}, \mathrm{target})$: claim batching
lives entirely on this reduction, via code that already exists. The
$\hat s_v$ folds are fused into the zerocheck and lincheck passes and
inherit their run-lists (an explicit checkpoint when generalizing those
passes), so the declared-work discipline extends here too.

\paragraph{The type converter: the virtual-opening sumcheck (new).} The
jagged transport consumes an \emph{evaluation} claim
$\hat f(\rho) = v$, not an inner product. Convert with the standard
reduction: a product sumcheck over the $M-7$ word variables for
$\sum_x f(x)\, b_{\mathrm{combined}}(x) = \mathrm{target}$. The prover's
work is support-proportional (both factors are materialized on declared
words only); at the final point $\rho$ the verifier evaluates
$\hat b_{\mathrm{combined}}(\rho)$ \emph{itself}, using the same
polylogarithmic machinery today's opening verifier uses at Ligerito's
residual point, and is left with the single evaluation claim
$\hat f(\rho) = v_\rho$. Everything non-factoring stays on the
verifier-computable side; the committed side crosses as one point claim.
Consequently \textbf{the jagged reduction always carries exactly one
claim per proof}. Historically this step is not new machinery at all: it
is \emph{ring-switching's own sumcheck}, which the direct path had fused
into Ligerito (itself an inner-product argument, able to absorb
$\langle f, b\rangle$ without an explicit sumcheck) and which the jagged
path un-fuses, because Ligerito now lives on the far side of the
dense/virtual index translation and only eq-shaped claims cross that
translation cheaply. The prover's sumcheck folds are support-proportional
under the declared-work discipline (implemented: on sparse supports the
$f$-side folds and the round messages touch only the live word
intervals), with one structural asymmetry worth recording: the $b$-side
folds are irreducibly \emph{dense}, for two stacked reasons. No product
structure survives ring-switching's per-word bit surgery, so folded
$b$-values cannot be produced on demand from small tables; and
support-\emph{skipping} fails as well: the late rounds consume folded
$b$ wherever folded $f$ is nonzero, and $f$'s support coarsens to the
\emph{full} domain as the folds merge dummy rows into live blocks ---
so, by backward induction (each round's folded array is the next
round's input), every entry of every round's $b$-array is load-bearing.
Skipping is sound only for a factor whose honest dead-region values are
\emph{zero} at every fold depth --- the $f$-side --- because zero is
the one value that need not be computed. What \emph{is} implemented
(sparse-gated): the full-domain $b$-array is never materialized ---
computing values on demand stores nothing, so nothing can go stale.
The round-0 message needs $b$ only at live pairs (dead pairs carry zero
witness factors) and is computed from the per-claim closed form there;
the round-0 fold then produces the half-size folded array directly from
the same closed form, so the $b$-side chain starts at half size and the
full-domain array never exists. The floor this leaves is intrinsic ---
one closed-form evaluation per padded word, since the late messages
depend on all of $b$ --- so the residual padded-domain cost is
compute, not memory. The verifier's costs --- the
static round count and one polylog $\hat b(\rho)$ evaluation --- are
oblivious to the sparsity pattern either way.

\paragraph{Reduction 2: the jagged transport.} Stack only the
\emph{declared} rows of the \emph{used} chunk-columns of all tables
contiguously (column-major per table --- the \code{BatchMajor} witness
layout, added in anticipation of exactly this), dropping the dummy rows,
the useless chunks, and the inter-slot gaps. This is \emph{height-$n_t$
stacking}, as implemented: slot $t$ contributes $\lceil u_t/128\rceil$
chunk-columns of $n_t$ words each --- the declared count, an arbitrary
integer in $[0, 2^{\nu}]$; non-power-of-two counts are native --- so the
un-padded dense size is the count-proportional
$\sum_t n_t\,\lceil u_t/128\rceil$ packed words (stacking granularity is
the $128$-variable chunk-column; the sub-chunk remainder is negligible).
The \emph{dense domain} --- the variable count of $\hat q$, hence the
round structure and the security configuration --- is this rounded up
to a power of two, \textbf{floored} at the smallest embedded Ligerito
security configuration ($m22$: $2^{15}$ packed words --- configurations
are derived and embedded per committed size) and capped at the union's
own padded size, so the dense domain never exceeds the virtual one. The
\emph{committed data}, by contrast, is no longer rounded: only an
integer number $L$ of Ligerito lanes covering the dense stack is
encoded and hashed, so the physical padding is below one lane (the
integer-lane commitment, below). The
jagged heights are therefore \textbf{per-proof} --- $n_t$ on slot $t$'s
used columns, $0$ elsewhere --- derived by prover and verifier alike
from the \emph{public} counts, never from the bundle; a wrong declared
count thus diverges in the commitment layout itself, in addition to the
transcript binding (\S\ref{sec:statement}) and the lincheck's
count-derived pin target (Remark~\ref{rem:pinbind}). Because the
union's virtual layout is itself BatchMajor (\S\ref{sec:union}), the
committed stack is simply the union buffer with dummy rows, useless
chunks, and gaps omitted ---
commitment and address space share one orientation, with no transpose
anywhere in the pipeline. Under the uniform-capacity convention the
identification is exact: a packed word's address is $[\,\nu \text{ batch
bits} \mid \text{column bits}\,]$ uniformly across slots, chunk-columns
enumerate in address order, and $\hat f$ \emph{is} the jagged grid
polynomial ($n = \nu$, $k = M - 7 - \nu$, per-proof heights $n_t$ words
on slot $t$'s used chunk-columns, $0$ elsewhere). Commit the dense stack
with Ligerito. The transport itself is
\emph{fused into the opening}: since $\hat f(\rho) = \sum_e q[e]\,
W_\rho(e)$ with $W_\rho(e) = \eq(\rho[0{:}\nu], \mathrm{row}(e))\cdot
\eq(\rho[\nu{:}], \mathrm{col}(e))$ --- every dropped word of the padded
buffer is zero, dummy rows by the partial-driver contract of
Remark~\ref{rem:padding} (debug-asserted by the transport's compaction)
and useless columns and gaps by construction, which is what keeps the
fused-opening identity $\ip{q}{W_\rho} = \hat f(\rho)$ intact (the
deleted terms of $\hat f(\rho)$ were all zero) --- Ligerito, itself an
inner-product argument, proves $\langle q, W_\rho\rangle = v_\rho$
directly with $W_\rho$ as its basis; no separate jagged sumcheck runs.
The factorization of $W_\rho$ is also how the prover \emph{sources} the
basis: $\eq(\rho[0{:}\nu], \mathrm{row}(e))$ and
$\eq(\rho[\nu{:}], \mathrm{col}(e))$ are two small $\eq$ tables plus a
column-cursor walk over \code{col\_prefix\_sums} --- valid at arbitrary
per-column heights --- so the $2^{M_{\mathrm{dense}}}$-word basis
vector is never materialized. The fused claim/round-0 pass and the
first fold fill cache-resident windows on demand; only the half-size
\emph{folded} basis that the first fold emits ever exists as an array,
and the later rounds consume it exactly as before.
The verifier's one remaining obligation is the basis MLE at Ligerito's
final points: the recursion ends not at a point but at a prover vector
$yr$ of length $2^{yr\_log\_n}$, requiring
$\hat W_\rho(\mathrm{ris}\,\|\,\mathrm{bits}(y))$ for every $y$. These
are supplied and bound by a \emph{send-and-spot-check} reduction: the
prover sends the vector
$\tilde b[y] = \hat W_\rho(\mathrm{ris}\,\|\,\mathrm{bits}(y))$; the
verifier uses it in the final check, squeezes a fresh
$r_{\mathrm{extra}}$, locally evaluates the MLE of the received
$\tilde b$ at $r_{\mathrm{extra}}$, and requires equality with a
single, unmodified assist evaluation
$\hat f_t(\rho[0{:}\nu], \rho[\nu{:}], \mathrm{ris}\,\|\,r_{\mathrm{extra}})$
--- by Schwartz--Zippel, a wrong $\tilde b$ survives with probability
$yr\_log\_n/|\F|$. \code{col\_prefix\_sums}, derived \emph{per proof}
from the public counts (the jagged heights are count-dependent under
height-$n_t$ stacking), encodes the layout inside the branching-program
evaluator $\hat f_t$.
The sizing fact, sharpened twice by implementation: above the $m22$
config floor the committed size is \textbf{count-proportional up to
less than one lane of padding}. The power-of-two rounding tax of the
previous draft (at most $2\times$; measured $28\%$ for a
full-utilization BLAKE3$+$SHA-256 mix at $M = 30$) is gone with the
integer-lane commitment below, and with it the column-dense rounding
effect (column-dense registries like BLAKE3$+$SHA-256, at $121$ of
$128$ and $246$ of $256$ chunk-columns, rounding straight back to the
padded power of two at full utilization). The floor binds only when the
dense area falls below $2^{15}$ packed words.

One consequence worth recording: under height-$n_t$ stacking the jagged
verifier's branching-program layout data (heights,
\code{col\_prefix\_sums}) is per-proof, derived from the public counts
--- which is where the assist delegation, already in place, earns its
keep a second time.

\paragraph{The integer-lane commitment (``integer commit, power-of-two
open'').} Ligerito's codeword is an \emph{interleaved} Reed--Solomon
encoding: the message is split into lanes of a fixed power-of-two
length $D$ ($2^{17}$ words at $M = 30$), one NTT per lane, and a Merkle
leaf holds one position of every lane. The lane count is where the
power-of-two rounding tax lived: rounding the dense stack up to
$2^{k}$ lanes RS-encoded and Merkle-hashed the padding like real data
($28\%$ of the committed stack for the full-utilization
BLAKE3$+$SHA-256 mix at $M = 30$; $34\%$ on the CLI's $\nu = 7$ tier).
The shipped scheme instead commits an \emph{arbitrary integer}
$L = \lceil N_{\mathrm{com}} / D \rceil$ of lanes --- padding below one
lane --- while the \emph{opening} keeps every power-of-two structure:
the fold still runs over the $2^{k}$-lane grid
($k = \lceil \log_2 L \rceil$), with lanes $[L, 2^{k})$
\emph{definitionally zero} --- the prover skips them, and the verifier
zero-fills each opened $L$-wide leaf before the consistency check ---
so the fold rounds, the lane $\eq$-tensor, $W_\rho$, $\tilde b$, and
the assist are all structurally unchanged (in particular the fused
transport survives). Two choices are load-bearing. \emph{High-bit
lanes}: lane $\ell$ is the contiguous block
$q[\ell D, (\ell{+}1) D)$, not the usual bit-strided interleaving ---
the dense stack's zero tail is then \emph{whole} zero lanes, and the
relabelling is a rotation of the index variables, which every
multilinear downstream survives by rotating its evaluation point. (With
strided lanes and integer $L$, the lane index $e \bmod L$ is not a bit
field, and a div/mod-$L$ would enter the jagged branching program,
blowing its one-bit carry up to $\approx 2L$ states.) \emph{Addressing,
not data}: the lane variables sit in the \emph{high} index bits, so
Ligerito's lane-binding rounds pair adjacent \emph{blocks} of $D$ words
instead of adjacent elements --- the same challenges, order, and $\eq$
weights, applied at block granularity --- and neither $q$ nor the basis
is ever permuted; the fast interleaved NTT kernel is kept verbatim. One
soundness point: the transcript binds only the commitment \emph{root},
while the opening reads the leaf width and lane count from the
commitment's parameters, which therefore ride the proof
attacker-controlled --- the union verifier requires them to equal the
count-derived parameters (a rejection, not a panic). When the stack
fills every lane ($L = 2^{k}$ --- in particular the
identity-compaction single-slot path) the scheme is byte-identical to
the power-of-two commit, which is what keeps the single-type anchors
green.

\begin{remark}[The direct-commitment harness]\label{rem:harness}
At full utilization ($n_t = 2^{\nu_t^{\max}}$ for all $t$) the union buffer
is itself a dense power-of-two array, and it can be committed directly with
today's \code{pcs::commit}; per-slot claims then transport \emph{trivially}:
for $\rho \in \F^{m_t}$, let $x^\star \in \F^{M}$ be $\rho$ in coordinates
$[0, m_t)$ and the bits of $p_t$ in coordinates $[m_t, M)$; then
$\hat z(x^\star) = \hat z_t(\rho)$. (Proof: decompose $\hat z$ by slot; slot
$t$'s prefix factor is $\eq(p_t, p_t) = 1$; for $t' \ne t$ with
$m_{t'} \le m_t$ the factor is $\eq$ of two distinct Boolean strings $= 0$;
for $m_{t'} > m_t$, if the top $M - m_{t'}$ bits of $p_t$ equaled $p_{t'}$,
slot $t$'s address range would lie inside slot $t'$'s, contradicting
disjointness.) This path is \textbf{not a protocol mode} and does not
appear in the wire format. It exists as (i) the differential-testing
harness for the union PIOP before and after the jagged layer lands, and
(ii) a cheaply resurrectable fallback should the jagged composition stall
--- the PIOP is identical under both transports.
\end{remark}

\subsection{Capacity-free ring-switching: a merged reduction (design
sketch)}\label{sec:merged}

Status: \textbf{shipped} --- since wire v6 the Mixed protocol's opening
IS this merged transport (design review passed 2026-07-27). The
load-bearing algebra is test-pinned (the Frobenius identity, the
linearized-coefficient extraction, the batched assist against the
brute-force twisted evaluation), and the full pipeline round-trips with
a tamper matrix --- proof-field tampers and commitment-params tampers
(the params-vs-root binding) --- on real mixed instances at all three
commit configurations (integer lanes at full and partial utilization,
power of two); measured results are folded into the cost paragraphs
below. The jagged transport remains in-tree as the differential/%
regression oracle (the M6 byte-pinned fixtures and the merged-vs-jagged
A/B tests), not a wire mode.

\paragraph{The padded-domain floor, measured.} Ring-switching is the one
pipeline stage that still sees the capacity: its weight and its sumcheck
live in padded coordinates, so the prover pays one $\Phi$-evaluation
(the $16{\cdot}256$-byte-table fold) per \emph{padded} word per claim,
plus the $b$-side fold chain. On the real $M = 30$ mixed workload
($16384 + 16384$) proved on capacity tiers $\nu = 14/15/16$ (100\%,
50\%, 25\% utilization), the combine $+$ virtual-opening sumcheck
measure $12.3 / 23.0 / 44.7$\,ms --- linear in $2^{M-7}$ --- while
every stage downstream of the type converter is flat (fused Ligerito
$\approx 17$--$18$\,ms, $\tilde b$/assist $\approx 1.5$\,ms at every
tier). At deployment scale this floor \emph{is} the capacity tax; the
support-proportional schemes of \S\ref{sec:commit} can reach the floor
(never below it), because the late round messages provably depend on
every entry of the weight.

\paragraph{The merged reduction.} The witness is identically zero off
the committed stack, so the type converter's sum restricts exactly:
\[
  \sum_{x} f(x)\, b(x) \;=\; \sum_{d} q[d]\; W[d],
  \qquad
  W[d] \;=\; \sum_i (\gamma_i \Phi_i)\bigl(\eq_{\mathrm{row},i}[\mathrm{row}(d)]
  \cdot \eq_{\mathrm{col},i}[\mathrm{col}(d)]\bigr),
\]
with $(\mathrm{row}, \mathrm{col})$ the jagged unrank, and
$W[d] := 0$ for $d \ge \mathrm{area}$ --- the sumcheck runs over the
full dense cube, so the weight must be defined on the power-of-two
tail; zero is the right extension on both sides ($q$'s committed tail
is zero, asserted, so the summand vanishes there either way, and the
assist's branching program computes exactly this zero-extension via
its comparison state, so the verifier's $\widehat W(\rho)$ agrees with
the prover's table by construction). Running the
virtual-opening sumcheck over the \emph{dense} domain with weight $W$
merges the two reductions of this section into one: it is
simultaneously ring-switching's sumcheck and the dense/virtual
translation. Scope: this replaces the \emph{union} path's transport
($K = 2$ ring-switch claims, no packed-direct claims); single-type
statements keep the current machinery, and the proof struct changes
(the $\tilde b$/spot-check fields go, the scalar $V$ and the assist
rounds arrive) --- a wire-format bump for the Mixed flavor. Prover cost becomes count-proportional --- the existing
just-in-time factored-weight machinery (the \code{col\_prefix\_sums}
cursor walk) with a $\Phi$ applied per element, roughly $6$\,ms at this
workload \emph{at any tier}. The sumcheck's output is a plain dense
evaluation $\hat q(\rho)$, so Ligerito reverts to an ordinary
$\eq$-basis opening --- the fused $W_\rho$ basis, the $\tilde b$
vector, and the send-and-spot-check reduction all retire.

\paragraph{The one new obligation, and why existing tools fail.} The
verifier must evaluate $\widehat W(\rho)$ --- the MLE over dense
coordinates of the $\Phi$-\emph{twisted} jagged weight. Untwisted, this
is exactly the assist's object $\hat f_t$; twisted, every existing tool
breaks on the same fact: $\Phi$ is $\F_2$-linear but
$\Phi(c\,x) \ne c\,\Phi(x)$, so it passes through neither the branching
program's multiplicative transitions, nor bit-decomposition (bit-slices
of field products are not multilinear), nor send-and-spot-check
(checking the sent value needs an independent oracle for the same
twisted object).

\paragraph{The bilinear-tensor branching program.} The twist does
untangle, at the price of state width. Write each term as
$\alpha(e)\cdot\Phi(\beta(e))$ with $\alpha = \eq(\rho, e)$ carrying
every $\F$-side weight and $\beta$ the row/column product inside
$\Phi$. The map $(\alpha, \beta) \mapsto \alpha\,\Phi(\beta)$ is
$\F_2$-bilinear, so
\[
  \widehat W(\rho) \;=\; B\Bigl(\;\sum_e \alpha(e) \otimes \beta(e)\Bigr),
  \qquad B(a \otimes b) := a\cdot\Phi(b),
\]
and the tensor $T = \sum_e \alpha \otimes \beta \in
\F \otimes_{\F_2} \F$ evolves under a \emph{read-once linear} branching
program: each index bit contributes the transparent Kronecker map
$M_{u(0)} {\otimes} M_{v(0)} + M_{u(1)} {\otimes} M_{v(1)}$
(multiplication by a constant is $\F_2$-linear), and the jagged
comparison/carry structure composes with it exactly as in the existing
width-$4$ program --- the state is (carry, comparison) $\otimes$ a
$128 \times 128$ bit matrix, and $\Phi$ is applied \textbf{once, to the
final state, outside every multiplication}.

\paragraph{Delegation analysis.} The implemented assist delegates
$\hat f_t$ by a sumcheck over the $2(m{+}1)$ column-\emph{boundary}
bits, against a width-$4$ program whose proof size is
width-independent (rounds $=$ variables) while prover cost and the
verifier's final DP scale linearly with width. It does \emph{not}
extend to the twisted object, for a reason worth recording precisely: a
sumcheck can only bind a variable that enters the summand
$\F$-linearly. The tensor form is $\F$-linear in the $\alpha$ slot ---
$B((M_x \otimes I)\,T) = x\,B(T)$ --- so $\rho$-side variables bind
fine; but the row-side factors (and the boundary bits that steer the
subtraction carry) live in the $\beta$ slot, where
$B((I \otimes M_x)\,T) \ne x\,B(T)$.

\paragraph{The resolution: Frobenius decomposition.} The obstruction
dissolves under a change of basis for $\Phi$ itself. Over
$\F_{2^{128}}$ every $\F_2$-linear map is uniquely a \emph{linearized
polynomial} $\Phi(x) = \sum_{j=0}^{127} c_j\,x^{2^j}$, and Frobenius
$x \mapsto x^{2^j}$ is a field \emph{automorphism}: it commutes with
products, and at a Boolean selector
$\eq(z, b)^{2^j} = \eq(z^{2^j}, b)$. Applied per term of the twisted
sum, $\beta(e)^{2^j}$ is the \emph{same} eq-product structure at
Frobenius-powered points, so
\[
  \widehat W(\rho) \;=\; \sum_{j=0}^{127} c_j\;
    \hat f_t\bigl(z_{\mathrm{row}}^{2^j},\, z_{\mathrm{col}}^{2^j},\,
    \rho\bigr)
\]
--- an $\F$-combination of $128$ \emph{ordinary, untwisted} assist
statements, with $\rho$ untouched (it lives in the $\alpha$ slot,
outside $\Phi$). The $\beta$-slot problem never arises: no
$\Phi$-interior variable is ever bound, because $\Phi$ is decomposed
into automorphisms that pass \emph{through} the product structure. The
load-bearing identity is verified numerically against the brute-force
jagged MLE (\code{frobenius\_twist\_matches\_assist\_object}: random
non-power-of-two heights, zero columns, Frobenius powers including
$j = 40$).

The delegation is therefore the \emph{existing} assist, $128$-way
batched (times the claim count): one sumcheck over the same $2(m{+}1)$
boundary-bit variables for $\sum_j \gamma_j\, W_j(c,d)\,\hat g_j(c,d)$,
proof size unchanged ($\approx 1.7$\,KB --- rounds $=$ variables,
independent of the batch width). Estimated costs: the verifier derives
the $c_j$ once per proof from $\Phi$'s byte table via a registry-static
Moore-matrix inverse ($128^2$ multiplies), and evaluates the $128$
per-$j$ weights by Frobenius chains --- the boundary factor is
$j$-independent and computed once, and
$\eq(z_{\mathrm{col}}^{2^j}, y) = \eq(z_{\mathrm{col}}, y)^{2^j}$ is
one squaring per column per $j$ --- so the verifier stays
\textbf{column-linear at $\approx 300$ operations per column}
($\approx 6\times$ today's assist path): it scales to registries of
thousands of columns, answering the scaling objection that killed the
fat verifier at registry scale. Prover: $\approx 128\times$ today's
(tiny) assist per claim --- a few ms parallel; the per-statement passes
are independent (an earlier draft claimed the suffix rows are global
Frobenius powers of the $j = 0$ rows --- false: they entangle the
\emph{unpowered} $\rho$-side factors with the powered $z$-side ones ---
only the layer tables come from squaring chains). The protocol and
soundness details follow. The \emph{fat verifier} ---
evaluating the twisted object directly by the per-column tensor DP, no
new argument at all --- survives as a fallback for small registries
($\approx 25$K multiplies per column: fine at today's $\approx 370$,
not at thousands; and unlike the comb there is no preprocessing escape,
since \code{col\_prefix\_sums} is per-proof).

\paragraph{What it buys, per tier --- measured.} The design's central
claim is \textbf{measured}: on the m30 load, prove-time growth from the
right-sized tier to $4\times$ over-capacity is $+14.1$\,ms, against the
jagged path's $+40.7$\,ms --- the merged transport is near-capacity-%
free, and the $256$--$512$\,MB $b_{\mathrm{combined}}$ transient is
gone at every tier. The prototype's $+34$\,ms right-sized-tier constant
was retired by the polish pass (integer-lane commit with the lane-major
eq-basis open, the segmented-cursor weight build, the seeded fused
combine for the lone eq claim, the block-dispatched assist prover):
warm in-run A/B at the right-sized tier reads $124.9$ vs.\
$130.9$\,ms in the merged path's favor, cold-gated runs put it within
$\approx 5$--$10$\,ms of the shipped constant (the batched assist's
remaining fixed cost), and it wins outright at $2\times$ capacity and
above. Proof: slightly smaller (the $\tilde b$ vector and
spot-check go, one batched assist arrives at unchanged size).
Verifier: $\lesssim 1$\,ms added (measured within noise at test
scale).

\subsubsection*{The batched Frobenius assist, in detail}

\emph{Notation.} Claims $i = 1..K$ (the union path has $K = 2$: the
$\alpha$-batched AB claim and the C claim), each with word-level point
$z_i = (z_i^{\mathrm{r}}, z_i^{\mathrm{c}})$ (rows low, columns high)
and $\F_2$-linear fold map $\gamma_i\Phi_i$ ($\gamma$ baked). For
$j = 0..127$ write $z^{(j)}$ for the coordinate-wise $2^j$-th power.
$T_0 \le T_1 \le \cdots$ are the per-proof column prefix sums
(\code{col\_prefix\_sums}); the assist's summation variables are the
boundary-pair bit strings $(u, v)$, $|u| = |v| = m{+}1$;
$U_{i,j}(u,v) = \sum_y \eq(z_i^{\mathrm{c},(j)}, y)\,
\eq(\mathrm{bits}(T_{y-1}), u)\,\eq(\mathrm{bits}(T_y), v)$ is the
boundary-indicator weight and $\hat g_{i,j}(u,v)$ the width-$4$
branching-program value with layer tables from
$(z_i^{\mathrm{r},(j)}, \rho)$ --- both exactly the implemented assist
objects at powered points.

\emph{Lemma 1 (linearized form).} Every $\F_2$-linear
$L : \F_{2^{128}} \to \F_{2^{128}}$ is uniquely
$L(x) = \sum_{j=0}^{127} c_j\,x^{2^j}$. The $c_j$ are recovered from
$L$'s values on the bit basis by one multiplication with the inverse
Moore matrix of that basis --- the matrix is universal (registry- and
proof-independent, precomputed once); per proof this costs $128$
evaluations of $L$ (byte-table applications) plus a $128^2$-multiply
matvec per claim.

\emph{Lemma 2 (Frobenius transport).} For Boolean $b$,
$\eq(z, b)^{2^j} = \eq(z^{2^j}, b)$ (squaring is additive in
characteristic $2$ and fixes $\F_2$), and Frobenius is a ring
automorphism, so it distributes over the eq-products; hence for every
$j$,
$\sum_e \eq(\rho, e)\,\beta(e)^{2^j} =
 \hat f_t(z^{\mathrm{r},(j)}, z^{\mathrm{c},(j)}, \rho)$
with $\rho$ untouched. (Pinned numerically by
\code{frobenius\_twist\_matches\_assist\_object}.)

\emph{Lemma 3 (restriction).} Define $f := \mathrm{embed}(q)$ (zero
off-stack --- the same definitional move the current fused transport
makes). Then $\sum_e f(e)\,b_{\mathrm{combined}}(e) =
\sum_d q[d]\,W[d]$ term-by-term: off-stack terms vanish with $f$, and
in-stack terms match under the unrank bijection. The merged sumcheck
therefore proves the ring-switch statement for the committed $q$'s
embedding, exactly as today's transport does.

\emph{Protocol.} After ring-switching ($\hat s_v$, $\gamma$-batching,
targets --- unchanged): (1) the \textbf{merged sumcheck}, $m_d$ rounds
over the dense domain for $\sum_d q[d]\,W[d] =
\mathrm{target}_{\mathrm{combined}}$, ending at point $\rho$ with
running claim $c^\ast$; (2) the prover sends $V$, claimed to be
$\widehat W(\rho)$; (3) the \textbf{combined assist sumcheck},
$2(m{+}1)$ rounds for
\[
  \textstyle\sum_{(u,v)} H(u,v) = V,
  \qquad
  H(u,v) := \sum_{i,j} c_{i,j}\; U_{i,j}(u,v)\;\hat g_{i,j}(u,v),
\]
with \textbf{deterministic weights}: the $c_{i,j}$ are
transcript-determined (Lemma 1 applied to $\gamma_i\Phi_i$), so no
fresh batching randomness is needed --- only the combined scalar $V$
is ever used, and plain sumcheck soundness on the single summand $H$
suffices. $H$ has degree $\le 2$ per variable, so rounds and proof
size match today's assist. (4) At the assist's final point $\sigma$
the verifier computes $H(\sigma)$ itself: the boundary factor
$E_y(\sigma)$ once ($(i,j)$-independent), the base column weights
$\eq(z_i^{\mathrm{c}}, y)$ once per claim, then one squaring per
column per $j$ (Frobenius chains), and $K \cdot 128$ width-$4$ DPs of
$O(m)$ each. (5) The verifier checks the merged sumcheck's final
relation $c^\ast = \hat q(\rho) \cdot V$, with $\hat q(\rho)$
discharged by the ordinary $\eq$-basis Ligerito opening ---
the fused $W_\rho$ basis, $\tilde b$, and send-and-spot-check retire.

\emph{Soundness.} Two plain sumchecks (errors
$2 m_d / |\F|$ and $4(m{+}1)/|\F|$), the PCS opening, and Lemmas 1--3;
the $c_{i,j}$ and $U, \hat g$ shapes are deterministic functions of
the transcript, so no new random-oracle surface. The write-up
obligation remaining is mechanical: transcript labels and order, and
the composition statement with ring-switching's existing argument.

\emph{Costs.} Verifier: $\approx 2(m{+}1)$ extra rounds to replay,
$E_y$ pass $O(\mathrm{cols} \cdot m)$, Frobenius chains
$O(\mathrm{cols} \cdot 128 \cdot K)$, DPs $O(128 K m)$, Moore matvec
$O(128^2 K)$ --- order $10^5$ operations at today's registry,
column-linear with constant $\approx 600$ at $K = 2$. Prover: the
$K \cdot 128$ statements' suffix rows and round passes are independent
($\approx 50$M multiplies at today's shapes, $2$--$4$\,ms parallel);
layer tables for all $j$ come from squaring chains. Proof:
$\approx 1.7$\,KB (unchanged rounds) plus the scalar $V$.

\emph{Implementation notes.} Zero-height columns collapse in
\code{assist\_columns} as today; the per-proof $T$ enters only through
$E_y$ and the DP transitions, unchanged; the $\Phi_i$ tables are the
existing \code{build\_fold\_byte\_table} outputs, consumed here only
through Lemma 1; at $K$ claims the statement count is $128K$, still
one sumcheck.

% ===========================================================================
\section{Statement, transcript, wire format}\label{sec:statement}

\paragraph{Statement digest and transcript.} The \emph{registry digest}
(\code{flock-registry-v1}) --- format version, the uniform capacity $\nu$,
and the ordered type list $(\kappa_t, u_t, \text{pin}_t)$ with each base
block's matrices --- is one precomputed 32-byte value. Per proof, the
transcript absorbs, before any challenge is squeezed: the
\code{flock-mixed-v1} domain label, the registry digest, the counts vector
$(n_t)_t$, and the commitment root. New domain-separation labels
(\code{flock-mixed-v1} for the statement binding,
\code{flock-virtual-open-v0} for the virtual-opening sumcheck) join the
existing family; the union lincheck keeps \code{flock-lincheck-v0},
differing from the single-table transcript only through its inputs.
Counts affect scalar bookkeeping --- run lists, fold extents, the
lincheck's count-derived pin target (Remark~\ref{rem:pinbind}) --- and,
under height-$n_t$ stacking (\S\ref{sec:commit}), the committed size:
both sides derive the dense length, hence the Ligerito configuration
(selected per proof from the embedded per-size table), from the public
counts, never from the bundle. The PIOP's round structure and the
verifier's control flow otherwise stay registry-static.

\paragraph{What a multi-table proof attests.} Exactly this: \emph{the
commitment opens to tables with the declared counts, every declared row of
which satisfies its type's hash relation.} The core deliberately carries
no statement about the hashed values themselves --- no per-invocation
input/output binding. All dataflow --- between tables, and between tables
and the public statement --- is the job of a planned \emph{lookup/bus}
layer built on top of this core (\S\ref{sec:open}): the core's only
obligations toward it, both already met, are fixed column offsets per type
(the \code{ChainLayout} region data) and cheap per-slot claims into the
transport batch. Until the bus lands, mixed proofs are infrastructure and
benchmarks, not application statements; today's single-type statements
(hash chains, Merkle paths) are untouched and keep their existing
claim-level glue.

\paragraph{Capacity tiers.} If workloads vary wildly, the registry may
offer a few \emph{capacity tiers} (small/medium/large instantiations of the
same type list); selecting a tier selects a registry digest. Each doubling
of capacity costs one static sumcheck round; see \S\ref{sec:open}.

\paragraph{Wire format (as shipped).} \code{proof\_io} is at
\code{VERSION} 6, which switches the \emph{Mixed} flavor's payload to
the merged transport of \S\ref{sec:merged} (v5 introduced the flavor
with the jagged-transport payload; the R1CS and Chain flavors ---
\code{hash\_kind} included --- are unchanged throughout; versioning is
strict, so older files are rejected). A mixed
bundle carries: the registry \emph{tier id} --- a stable one-byte code
that pins the \emph{full} registry, base matrices, widths, and the uniform
capacity $\nu$ included (current tiers: BLAKE3$+$SHA-256 at $\nu = 7$ and
$\nu = 10$) --- the counts vector (one \code{u64} per type, in slot
order), the commitment, and the merged-transport union proof. CLI: \code{flock\_chain
prove --mix blake3=N,sha2=M} proves one mixed proof on the smallest
fitting tier; \code{verify} auto-detects the bundle flavor. The statement
remains well-formedness only (no per-invocation IO); the existing
single-hash chain statements keep their own flavor and path, untouched.

% ===========================================================================
\section{Cost accounting}\label{sec:costs}

Let $N_{\mathrm{decl}} = \sum_t n_t k_t$ (declared area),
$N_{\mathrm{com}} = \sum_t n_t \lceil u_t/128\rceil$ packed words (the
count-proportional dense size under height-$n_t$ stacking, committed
as an integer number of Ligerito lanes --- physical padding below one
lane --- with the opening's power-of-two domain floored at the $m22$
config, \S\ref{sec:commit}), $T$ = types.

\begin{center}
\begin{tabular}{@{}ll@{}}
\toprule
 & \textbf{cost} \\
\midrule
Witness generation & per-type closures over declared rows; nothing for dummies \\
Commit & Ligerito on $N_{\mathrm{com}}$ (declared rows of used chunk-columns --- count-proportional) \\
Zerocheck prover & $O(N_{\mathrm{decl}})$; static $M$ rounds; skip intact \\
Lincheck prover & $O(N_{\mathrm{decl}})$ folds $+$ $(M{-}\nu)$-round union-column SC \\
Ring-switch folds & $O(N_{\mathrm{decl}})$ (fused into the zerocheck/lincheck passes) \\
Virtual-opening SC prover & $O(\text{support})$, $(M{-}7)$-round product SC \\
Jagged transport (fused into Ligerito) & $\tilde O(N_{\mathrm{com}})$ (factored $W_\rho$ basis, just-in-time, $+$ $\tilde b$/assist) \\
Opening & one batched Ligerito open \\
\midrule
Verifier rounds & $M + K$ registry-static; transport rounds sized by the per-proof $N_{\mathrm{com}}$ \\
Verifier matrix work & $O(\sum_t \nnz(A_t,B_t))$ \\
Verifier bookkeeping & $O(T \cdot M)$ weights $+$ per-proof jagged layout data from the counts (assisted) \\
Proof size & $+\,O(T)$ scalars over a single-table proof of equal area \\
\bottomrule
\end{tabular}
\end{center}

The headline: \textbf{the prover pays for declared work in the PIOP} ---
witness generation and zerocheck folding match the single-table system on
the same workload; capacity costs only $M - \log_2 N_{\mathrm{decl}}$
extra sumcheck rounds, and mixing costs $O(T)$ scalars plus the (new, but
size-independent-per-invocation) jagged transport. The commit pays for
the declared rows of the used columns --- count-proportional above the
$m22$ config floor under height-$n_t$ stacking, with physical padding
below one Ligerito lane under the integer-lane commitment
(\S\ref{sec:commit}); the former power-of-two rounding tax, and the
column-dense effect it produced at full utilization, are gone.

% ===========================================================================
\section{Roadmap}\label{sec:roadmap}

Jagged is on the critical path for everything, so it lands first, in the
smallest possible context.

\begin{description}[leftmargin=1.6em, style=nextline]
  \item[Phase 0 --- enablers (no behavior change).]
    Registry/schedule types; \code{BlockR1cs} $\to$ slot schedule
    (single-table $=$ one slot; all existing tests must pass unchanged);
    run-list \code{PaddingSpec}; registry digest. Gate: existing
    single-hash proofs byte-stable (modulo version field) and benchmarks
    flat.
  \item[Phase 1 --- jagged under the existing single-table system.]
    Wire the three-polynomial pipeline (\S\ref{sec:commit}) beneath
    today's statements (one-type registry, the existing BatchMajor prover
    paths): keep ring-switching and $\gamma$-batching verbatim, add the
    virtual-opening sumcheck and the single-claim jagged transport with
    assist, and open the dense stack. Same statements, new commitment
    path, so every existing test and benchmark is a correctness and
    throughput oracle. Gate: single-hash throughput and proof size within
    an agreed margin of the direct path. \emph{No area saving is expected
    here} --- at single-table full utilization the power-of-two committed
    length rounds the dense stack back to the padded size
    (\S\ref{sec:commit}); the saving is a Phase 2 gate.
  \item[Phase 2 --- the union PIOP and dynamic counts (mixed proofs ship).]
    Union zerocheck (count-derived run lists); union-column lincheck
    (\S\ref{sec:multisize}); per-slot witness
    closures in the existing drivers; counts in the
    transcript; wire format v5; CLI \code{--mix}. Internal milestone: the
    union PIOP differentially tested at full utilization against the
    direct-commitment harness (Remark~\ref{rem:harness}). Gates: mixed
    BLAKE3$+$SHA-256 proof verifies end-to-end; prover time tracks
    \emph{declared} work across utilization levels; verifier time flat
    across utilization levels; throughput within noise of the weighted
    average of the single-hash provers on the same workload;
    count-proportional committed-area saving delivered by height-$n_t$
    stacking (\S\ref{sec:commit}) --- the measured gate: BLAKE3$+$SHA-256
    at $\nu = 7$, counts $(32, 32)$, commits $2^{15}$ words (the $m22$
    config floor) against capacity-height stacking's $2^{16}$; at
    $\nu = 10$ low utilization the saving reaches $16\times$ ($2^{15}$
    vs $2^{19}$ committed words).
\end{description}

Independent invocations keep every phase per-table testable: any mixed
proof can be differentially tested against the corresponding single-table
proofs on the same witnesses.

\subsection*{Implementation status}

Phases 0--2 are implemented on this branch, under the oracle discipline
the roadmap prescribes: every step landed behind a single-slot
byte-identity differential against the direct path (the whole proof
bundle asserted byte-identical, not merely accepted), with brute-force
MLE cross-checks for the $T = 2$ lincheck; the direct-commitment harness
(Remark~\ref{rem:harness}) remains a live regression anchor, and the
optimization work below additionally landed behind SHA-256-pinned
proof-byte fixtures at four count vectors (full, partial,
non-power-of-two, and zero counts) --- a permanent regression anchor
for the mixed path. Height-$n_t$
stacking (M5) made the committed size
count-proportional (\S\ref{sec:commit}) with the area numbers of the
Phase 2 gate above: BLAKE3$+$SHA-256 at $\nu = 7$, counts $(32, 32)$,
commits $2^{15}$ words against capacity-height's $2^{16}$, and $16\times$
fewer at $\nu = 10$ low utilization. Since draft v0.4:

\begin{itemize}[nosep]
  \item \emph{Declared-work prover passes}
        (item~\ref{item:folding} delivered): the zerocheck tail rounds,
        the union lincheck row folds (count-proportional at every
        partial count), and the virtual-opening $f$-side folds are
        support-proportional; $\tilde b$ is one batched shared-prefix
        branching-program evaluation; and the run-list zerocheck path
        shares one kernel with the single-run path (a missed-SIMD gap
        measured at $+13\%$ zerocheck on identical data, now
        $+0.5\%$). At $\nu = 10$ counts $(8, 8)$ the prove dropped
        $\approx 1.34\times$ (single-shot wall $\approx 1.5\times$);
        full utilization is unchanged or slightly better. The residual
        low-utilization floor is count-independent by nature: per-type
        comb builds ($O(\nnz)$, registry-static), the $m22$-floor
        commit/opening, and the dense $b$-side (\S\ref{sec:commit}).
  \item The \emph{integer-lane commitment} (\S\ref{sec:commit}): at
        $M = 30$ full-utilization BLAKE3$+$SHA-256 the commit encodes
        $46$ of the former $64$ lanes (NTT $22.6 \to 15.5$\,ms, Merkle
        $8.9 \to 6.7$\,ms) and the proof shrinks $16\%$
        ($400\,768 \to 336\,288$ bytes); the CLI's
        \code{blake3=100,sha2=37} proof shrinks $28\%$
        ($274 \to 198$\,KB). With the lane binding done by addressing
        rather than data movement, prove is $\approx 6\%$ \emph{faster}
        than the power-of-two-commit baseline, on top of the size win.
  \item The \emph{just-in-time factored basis} (\S\ref{sec:commit}):
        the fused opening no longer materializes $W_\rho$ ($134$\,MB at
        $M = 30$) --- the fused open dropped $16.0 \to 14.4$\,ms and
        peak memory by the full $134$\,MB.
  \item (v0.7) The \emph{merged jagged/ring-switch transport}
        (\S\ref{sec:merged}) shipped as the Mixed flavor's opening ---
        wire \code{VERSION} 6; design review passed 2026-07-27. The
        jagged transport and its byte-pinned fixtures stay in-tree as
        the differential oracle.
\end{itemize}

Measured at $m = 30$ BLAKE3 on Apple silicon: the direct path proves at
$\approx 600$k compressions/s ($\approx 108$\,ms); the union pipeline
runs at $\approx 96\%$ of direct, the remaining structural cost being
essentially the virtual-opening sumcheck ($\approx 3.8$\,ms --- it can
never fuse, \S\ref{sec:commit}). A mixed BLAKE3$+$SHA-256 proof at the
same scale costs $\approx 1.1\times$ the sum of the corresponding
single-type proofs ($\approx 300$k invocations/s) at $336$\,KB; the
CLI's mixed proof of \code{blake3=100,sha2=37} is $\approx 198$\,KB,
proving and verifying in $\approx 0.3$\,s each. Verification is
$\approx 11$\,ms, of which $\approx 83\%$ is the per-type comb builds
--- a cost the direct path pays too, fused into its sumcheck replay
(item~\ref{item:manytypes}).

The merged jagged/ring-switch transport of \S\ref{sec:merged} is
\textbf{shipped as of wire v6} (\code{open\_batch\_merged}, the
Frobenius assist, the \code{\_merged} union prove/verify pair behind
the Mixed flavor): end-to-end roundtrip and tamper matrix --- proof
fields and commitment params --- on real mixed instances,
capacity-freedom measured at the m30 load ($+14.1$ vs $+40.7$\,ms
growth to the $4\times$ tier); the jagged transport stays in-tree as
the differential oracle (\S\ref{sec:merged}).

Open follow-ups, in the order they are likely to matter: a Keccak
batch-major driver (the registry machinery accepts the type; only the
driver is missing); the lookup/bus layer (\S\ref{sec:open}); and
comb delegation (item~\ref{item:manytypes}), explicitly deferred until
registry growth or verify latency demands it.

% ===========================================================================
\section{Open questions}\label{sec:open}

\begin{enumerate}[nosep]
  \item \textbf{Phase 1 wiring constraints (largely resolved).} Claim
        batching is fully resolved by architecture (\S\ref{sec:commit}:
        it lives on reduction 1, in existing $\gamma$-machinery; the
        transport always carries one claim). The Ligerito-configuration
        question is settled by per-size enumeration: under height-$n_t$
        stacking the dense size (\code{dense\_m}) is count-dependent, so
        security configurations are derived and embedded per committed
        size and the configuration is selected \emph{per proof} from the
        public counts, floored at the smallest embedded configuration
        ($m22$, \S\ref{sec:commit}). The virtual-opening
        sumcheck's domain label (\code{flock-virtual-open-v0}) and
        transcript position are pinned. The last residue --- the
        ring-switch \code{ChunkPadding} multi-run fallback --- turned
        out to be moot on the union path: the union prover hands the
        opening \emph{precomputed} slice-folds $\hat s_v$, captured
        during its own PIOP passes, so the ring-switch fold the
        fallback guards never runs there. Nothing in this item remains
        open.
  \item \textbf{Capacity sizing.} How many tiers, and how large? Each
        doubling of capacity costs one static sumcheck round and slightly
        larger bookkeeping tables; exceeding capacity forces a proof
        split. Leaning: a single generous uniform-capacity tier
        (\S\ref{sec:model}), adding tiers only on demonstrated demand;
        confirm against a workload model at Phase 2.
  \item \label{item:folding}\textbf{Folding over empty capacity --- an
        implementation
        pattern, not a hot loop.} The union fold factorizes per slot:
        binding variables low to high, slot $t$ folds its own dense,
        declared data for its first $m_t$ rounds (today's kernels,
        verbatim) and then collapses to a single scalar; in every later
        round it contributes $O(1)$ scalar work (the accumulated value
        times the appropriate $\eq$ factors --- the implicit-batching
        identity \eqref{eq:implicit-batch}, operationally). So ``prover
        pays declared work'' is delivered by architecture --- per-slot
        dense sub-folds plus an $O(T)$-per-round scalar tail --- with no
        sparse iteration inside any hot loop. The same discipline covers
        the ring-switching folds (\S\ref{sec:commit}): the fused
        $\hat s_v$ capture and the per-slot claim folds must skip dummy
        rows via the same run-lists. \emph{Delivered}: implemented
        end-to-end behind proof-byte fixtures, with the adversarial
        low-utilization benchmark in place (Implementation status,
        \S\ref{sec:roadmap}); the one structural exception is the
        $b$-side of the virtual-opening sumcheck, which is irreducibly
        dense (\S\ref{sec:commit}).
  \item \label{item:manytypes}\textbf{Many types.} Verifier matrix work
        is linear in registry
        size ($\sum_t \nnz$). If the registry grows past a handful of
        types, the clean fix is Spartan-style preprocessing: commit each
        base block's MLE once at registry-definition time and replace the
        verifier's comb evaluation with one PCS opening per type;
        assist-style delegation is the lighter-weight alternative. Now
        measured: the comb evaluation is $\approx 83\%$ of the
        $\approx 11$\,ms verify at $m = 30$. The direct path pays the
        same cost, fused into its sumcheck replay --- this is the
        lincheck's verify cost, not a multi-table artifact; a
        multi-slot registry multiplies it by the number of
        \emph{distinct} circuits. Delegation is explicitly deferred
        until registry growth or verify latency demands it --- and it
        pairs naturally with the merged reduction of \S\ref{sec:merged},
        whose leading variant spends verifier time to delete the
        prover's padded-domain floor.
  \item \textbf{The padded-domain floor / merged reduction.} The one
        remaining capacity-proportional prover cost is ring-switching's
        padded-domain $\Phi$-pass; \S\ref{sec:merged} sketches the
        merged jagged/ring-switch reduction that would make it
        count-proportional. The formerly open $\Phi$-twisted weight
        evaluation is resolved in principle by the Frobenius
        decomposition (a $128$-way batched ordinary assist; core
        identity test-verified); the remaining work is the batched
        assist protocol and its soundness write-up.
  \item \textbf{The lookup/bus layer (the designated follow-on).} All
        dataflow --- table-to-table communication \emph{and} binding
        hashed values to the public statement --- is to be provided by a
        lookup/bus argument over the committed columns: rows emit and
        consume tuples (e.g.\ (type, input, output)); a multiset-balance
        argument (grand-product or log-derivative/LogUp style, realized as
        additional sumchecks over committed columns) enforces that every
        consumed tuple was emitted; public values enter as statement-side
        injections into the bus. This integrates at claim level --- the
        bus sumchecks emit evaluation claims into the same transport batch
        --- so the multi-table core needs nothing from it beyond what
        already holds (fixed column offsets per type, cheap per-slot
        claims). It changes no decision in this document and deserves its
        own design document.
  \item \textbf{Internal-hash generality.} Orthogonal axis (Merkle/FS hash
        is hard-coded SHA-256); out of scope here, noted for completeness.
\end{enumerate}

% ===========================================================================
\appendix
\section{Code touchpoint map}\label{app:code}

The design-time map, kept for orientation; the changes have since landed
--- the registry and slot schedule in \code{flock-core/src/schedule.rs},
the union layer in \code{union.rs}, the union lincheck in
\code{lincheck/union.rs}, and the tier table in
\code{flock-prover/src/mixed.rs}.

\begin{center}
\small
\begin{tabular}{@{}lll@{}}
\toprule
\textbf{Area} & \textbf{Location (today)} & \textbf{Change} \\
\midrule
Instance & \code{flock-core/src/r1cs.rs} (\code{BlockR1cs}) & slot schedule; \code{apply\_a} per slot \\
Zerocheck & \code{flock-core/src/zerocheck.rs} (\code{PaddingSpec}) & count-derived run lists; kernel unchanged \\
Lincheck & \code{flock-core/src/lincheck.rs} & multi-size inner SC; prefix weights; \\
 & (\code{LincheckCircuit}, \code{fold\_alpha\_batched}) & per-type comb reuse \\
Commitment & \code{flock-core/src/pcs/jagged.rs} $+$ branch & wire in as the claim-transport layer \\
Claim batching & \code{flock-core/src/pcs.rs} (\code{open\_batch\ldots}) & feeds the jagged transport batch \\
Test harness & \code{flock-core/src/pcs/commit.rs} & direct commit at full utilization \\
 & & (Remark~\ref{rem:harness}; tests only) \\
Statement & \code{flock-core/src/r1cs.rs} (\code{statement\_digest}) & registry digest $+$ counts \\
Witness gen & \code{flock-prover/src/r1cs\_hashes/common.rs} & drivers over per-slot closures \\
Statement glue & \code{flock-prover/src/r1cs\_hashes/chain\_common.rs} & unchanged (single-type statements); \\
 & & superseded later by the bus layer \\
Wire/CLI & \code{flock-prover/src/proof\_io.rs}, \code{bin/flock\_chain.rs} (retired 2026-08-14) & v5; registry id $+$ counts; \code{--mix} \\
\bottomrule
\end{tabular}
\end{center}

\end{document}
